\section{A fermion forces a prescribed Morita self-duality}
\label{sec:fermion-azumaya-crossing}

The rank-$22$ based module in the preceding section must satisfy an
additional identity that uses the braiding of the rank-$26$ candidate.
This condition prescribes its Morita dual, instead of allowing an
arbitrary dual Gram factorization.

\begin{lemma}[Fermion Azumaya crossing]
\label{lem:fermion-azumaya-crossing}
Let $\mathcal B$ be a braided fusion category containing a fermion $f$,
and let $A=\mathbf1\oplus f$ be its group algebra as an algebra object.
For simple right $A$-modules $m_a$, put
\[
 M_i[a,b]=\dim\operatorname{Hom}_{\mathcal B_A}(i\otimes m_a,m_b),
 \qquad i\in\operatorname{Irr}(\mathcal B).
\]
Then, for all module labels $a,b,c,d$,
\begin{equation}
 \sum_i M_i[a,b]M_i[c,d]
 =\sum_i M_i[c,a]M_i[d,b].
 \label{eq:fermion-azumaya-crossing}
\end{equation}
\end{lemma}
\begin{proof}
Inside $\langle f\rangle\simeq\mathrm{sVec}$, the algebra $A$ is the
one-generator Clifford algebra.  Its Azumaya map
$A\otimes A^{\mathrm{op}}\longrightarrow\underline{\operatorname{End}}(A)$
is an isomorphism: the identity, left multiplication by the odd generator,
graded right multiplication by that generator, and their composite form
a basis of the four-dimensional endomorphism algebra.  The braided
embedding into $\mathcal B$ preserves this isomorphism.  Thus $A$ is
Azumaya, and both alpha-induction functors
$\mathcal B\longrightarrow {}_A\mathcal B_A$ are tensor equivalences.
Here we use the Azumaya criterion in Davydov--Nikshych,
\emph{The Picard crossed module of a braided tensor category},
Section~3.2, Definition~3.5 and the paragraph following it
(\href{https://arxiv.org/abs/1202.0061}{arXiv:1202.0061}).

Under either equivalence, the right action on $\mathcal B_A$ has the
same multiplicity matrices as the left action; the braiding identifies
the corresponding products.  Associativity in the Morita linking
category, applied to $(m_a m_b^*)m_c=m_a(m_b^*m_c)$, therefore gives
\[
 \sum_i M_i[b,a]M_i[c,d]
 =\sum_i M_i[b,c]M_i[a,d].
\]
Exchange $a,b$, and use $M_{i^*}=M_i^{\mathsf T}$ while reindexing the
right-hand sum by $i\mapsto i^*$.  This proves
\eqref{eq:fermion-azumaya-crossing}.
\end{proof}

\paragraph{Application to $\RR{4}{155}$.}
For the retained rank-$26$ modular candidate, reconstruct the $26$
action matrices from the previously recorded six-dimensional
anti-invariant witness.  All $22^4=234256$ identities
\eqref{eq:fermion-azumaya-crossing} hold exactly.  The replay uses integer
arithmetic and binds both frozen inputs by SHA256; it is recorded in
\filename{evidence/unified_20261016/rank6/verify_fermion_crossing.py}.
Thus the known based module is compatible with the prescribed Azumaya
Morita self-duality.  The remaining problem is realization by coherent
module associators.  This passing identity constructs no module category
and proves no categorification or exclusion; both rank-six cases remain
unresolved.  Braidedness is essential in the lemma: commutativity of a
Grothendieck ring alone does not identify the two Morita actions.
